\documentclass[border=4pt]{standalone}
\usepackage[europeanresistor, europeaninductor]{circuitikz}
\usepackage{siunitx}
\usepackage{textcomp}
\ctikzset{bipoles/length=1.4cm}
\begin{document}
\begin{circuitikz}
% ===== left cluster: (AC source + 1 ohm) parallel with 1 mOhm, selected by SPDT =====
\draw (2,6.4) -- (4,6.4);                          % top wire joining the two branches
\draw (3,6.4) node[circ]{};                         % junction
\draw (3,6.4) -- (3,7.0) -- (7,7.0) -- (7,6.4);     % up and over to the right branch
% source branch: AC source in series with 1 ohm
\draw (2,6.4) to[sV, l_={\qty{1}{\volt}/\qty{20}{\kilo\hertz}}] (2,5.0);
\draw (2,5.0) to[R, l_={\qty{1}{\milli\ohm}}] (2,3.6);
% parallel branch: 1 mOhm (auto-centred over the taller span)
\draw (4,6.4) to[R, l={\qty{1}{\milli\ohm}}] (4,3.6);
% two-way (changeover) switch: common at bottom, throws up to each branch
\node[spdt, rotate=90] (Sw) at (3,2.9) {};
\draw (Sw.out 1) |- (2,3.6);
\draw (Sw.out 2) |- (4,3.6);
\draw (Sw.in) -- (3,1.4);                           % common down to the capacitor
% ===== DC-link capacitor (vertical, below the switch) =====
\draw (3,1.4) to[C, l_={\qty{0.37}{\milli\farad}}] (3,0);
% ===== right branch: ammeter + field ROM (stacked) =====
\draw (7,6.4) to[rmeter, t=A] (7,5.0);
\draw (7,5.0) -- (7,4.6);
\draw (7,4.6) to[sV, l={field ROM}] (7,3.2);
\draw (7,3.2) -- (7,0);
% ===== bottom rail: right branch -> ground -> capacitor =====
\draw (7,0) -- (3,0);
\draw (5,0) node[ground]{};
\end{circuitikz}
\end{document}